\chapter{Green-Relationen und innere Struktur}
\label{b28guide:chap:green}

Die Green-Relationen und die folgenden Klassenaussagen sind für jede
Halbgruppe definiert. Die vollständige Alternative für eine
$\mathcal H$-Klasse und die Schützenbergergruppe behandeln jede solche
Klasse, ohne besondere Elemente oder Gruppenklassen vorauszusetzen.

\section{Hauptideale und Green-Relationen}

\begin{definition}[Green-Präordnungen und Green-Relationen]
\label{b28guide:def:green}
Für $a,b\in S$ seien
\[
\begin{aligned}
a\leq_{\mathcal R}b&\iff a\in bS^1,&
a\leq_{\mathcal L}b&\iff a\in S^1b,\\
a\leq_{\mathcal J}b&\iff a\in S^1bS^1.
\end{aligned}
\]
Für $\mathcal K\in\{\mathcal R,\mathcal L,\mathcal J\}$ bedeutet
$a\mathrel{\mathcal K}b$, daß sowohl $a\leq_{\mathcal K}b$ als auch
$b\leq_{\mathcal K}a$ gelten. Weiter sind
\[
\mathcal H=\mathcal R\cap\mathcal L,
\qquad
\mathcal D=\text{kleinste Äquivalenzrelation über }
\mathcal R\cup\mathcal L.
\]
Die Klassen von $a$ heißen $R_a,L_a,J_a,H_a,D_a$.
\end{definition}

\begin{theorem}[Teilbarkeit und Inklusion der Hauptideale]
\label{b28guide:thm:green-preorders}
Die drei Green-Präordnungen sind reflexiv und transitiv. Für $a,b\in S$
gilt
\[
\begin{aligned}
a\leq_{\mathcal R}b&\iff aS^1\subseteq bS^1,\\
a\leq_{\mathcal L}b&\iff S^1a\subseteq S^1b,\\
a\leq_{\mathcal J}b&\iff S^1aS^1\subseteq S^1bS^1.
\end{aligned}
\]
Die Relationen $\mathcal R,\mathcal L,\mathcal J,\mathcal H,\mathcal D$
sind Äquivalenzrelationen. Insbesondere gelten
\[
\begin{aligned}
a\mathrel{\mathcal R}b&\iff aS^1=bS^1,\\
a\mathrel{\mathcal L}b&\iff S^1a=S^1b,\\
a\mathrel{\mathcal J}b&\iff S^1aS^1=S^1bS^1.
\end{aligned}
\]
Auf den $\mathcal K$-Klassen definiert
$K_a\leq K_b\iff a\leq_{\mathcal K}b$ jeweils eine partielle Ordnung
für $\mathcal K\in\{\mathcal R,\mathcal L,\mathcal J\}$.
\cite{howie,green1951}
\end{theorem}

\begin{definition}[Untere Mengen von Hauptklassen]
Eine Teilmenge $\mathcal A\subseteq S/\mathcal J$ heißt \emph{untere Menge},
wenn aus $J\in\mathcal A$ und $J'\leq J$ stets $J'\in\mathcal A$ folgt.
Die Ordnung auf $S/\mathcal J$ ist dabei die oben definierte
Inklusionsordnung der zweiseitigen Hauptideale.
\end{definition}

\begin{theorem}[Ideale als untere Mengen der $\mathcal J$-Klassen]
\label{b28guide:thm:ideals-j-order}
Die Abbildungen
\[
 I\longmapsto\{J_a:a\in I\},
 \qquad
 \mathcal A\longmapsto\bigcup_{J\in\mathcal A}J
\]
sind zueinander inverse vollständige Verbandsisomorphismen
zwischen den Idealen von $S$ einschließlich $\varnothing$ und
den unteren Mengen von $S/\mathcal J$.
Insbesondere ist jedes Ideal eine Vereinigung ganzer $\mathcal J$-Klassen;
die geordnete Menge dieser Klassen bestimmt den gesamten Idealverband.
\cite{koussas2021,howie}
\end{theorem}
\begin{theorem}[Verträglichkeit mit einseitiger Multiplikation]
\label{b28guide:thm:green-compatibility}
Für $a,b\in S$ und $s\in S^1$ gelten
\[
\begin{aligned}
a\leq_{\mathcal R}b&\Longrightarrow sa\leq_{\mathcal R}sb,&
a\mathrel{\mathcal R}b&\Longrightarrow sa\mathrel{\mathcal R}sb,\\
a\leq_{\mathcal L}b&\Longrightarrow as\leq_{\mathcal L}bs,&
a\mathrel{\mathcal L}b&\Longrightarrow as\mathrel{\mathcal L}bs.
\end{aligned}
\]
Somit ist $\mathcal R$ eine Linkskongruenz und $\mathcal L$ eine
Rechtskongruenz. Für beliebige $a,b\in S$ gilt außerdem
\[
ab\leq_{\mathcal R}a,\qquad ab\leq_{\mathcal L}b,
\qquad ab\leq_{\mathcal J}a,\qquad ab\leq_{\mathcal J}b.
\]
\cite{howie}
\end{theorem}

\begin{theorem}[Verknüpfung der Green-Relationen]
\label{b28guide:thm:green-d}
Es gelten
\[
\mathcal H\subseteq\mathcal R,\mathcal L
\subseteq\mathcal D\subseteq\mathcal J,
\qquad
\mathcal D=\mathcal R\circ\mathcal L
           =\mathcal L\circ\mathcal R.
\]
Ausgeschrieben bedeutet dies für $a,b\in S$:
\[
\begin{aligned}
a\mathrel{\mathcal D}b
&\iff (\exists c\in S)\;
       a\mathrel{\mathcal R}c\mathrel{\mathcal L}b\\
&\iff (\exists d\in S)\;
       a\mathrel{\mathcal L}d\mathrel{\mathcal R}b.
\end{aligned}
\]
Eine $\mathcal J$-Klasse ist eine Vereinigung von $\mathcal D$-Klassen;
die Gleichheit $\mathcal D=\mathcal J$ gilt nicht für jede Halbgruppe.
\cite{howie,green1951}
\end{theorem}

\section{Mitschs natürliche partielle Ordnung}

\begin{definition}[Natürliche Ordnung]
Für $a,b\in S$ definieren wir
\[
\begin{aligned}
 a\leq_M b\iff{}&a=b\ \text{oder}\\
 & (\exists x,y\in S)\;a=xb=by\ \text{und}\ xa=a.
\end{aligned}
\]
Dies definiert die Vergleichbarkeit einzelner Paare; eine Einheit von $S$
oder idempotente oder invertierbare Faktoren werden nicht vorausgesetzt.
\end{definition}

\begin{theorem}[Allgemeiner Ordnungssatz von Mitsch]
\label{b28guide:thm:mitsch}
Für jede Halbgruppe $S$ ist $\leq_M$ eine partielle Ordnung:
Sie ist reflexiv, antisymmetrisch und transitiv.
Mit der stets möglichen Einheitserweiterung $S^1$ aus
Definition~\ref{b28guide:def:unit} lautet die gleichwertige Kurzform
\[
 a\leq_M b
 \iff
 (\exists x,y\in S^1)\;
 a=xb=by,\quad xa=a=ay.
\]
Ferner gilt für alle $a,b\in S$:
\[
 a\leq_M b
 \ \Longrightarrow\
 a\leq_{\mathcal R}b\ \text{und}\ a\leq_{\mathcal L}b.
\]
Eine Verträglichkeit der Ordnung mit beliebiger Links- oder
Rechtsmultiplikation ist im allgemeinen nicht gegeben.
\cite{mitsch1994}
\end{theorem}

\begin{theorem}[Ordnungssatz von P. R. Jones]
\label{b28guide:thm:jones-order}
Für $a\in S$ sei
\[
 \mathord{\downarrow}_M a=\{x\in S:x\leq_M a\},
\]
versehen mit der Einschränkung der Mitsch-Ordnung.
Für jede Halbgruppe $S$ und alle $a,b\in S$ mit
$a\mathrel{\mathcal D}b$ sind die geordneten Mengen
$\mathord{\downarrow}_M a$ und $\mathord{\downarrow}_M b$
ordnungsisomorph. Es gibt also eine Bijektion
$\varphi:\mathord{\downarrow}_M a\to\mathord{\downarrow}_M b$ mit
\[
 x\leq_M y\iff\varphi(x)\leq_M\varphi(y)
 \qquad(x,y\in\mathord{\downarrow}_M a).
\]
Der Ordnungstyp hängt somit nur von der $\mathcal D$-Klasse ab.
Idempotenz oder Regularität der Elemente oder von $S$ werden
nicht vorausgesetzt. \cite{mitsch1994}
\end{theorem}

\section{Greens Lemma und das Klassenraster}

\begin{theorem}[Greens Lemma: Rechtsmultiplikation]
\label{b28guide:thm:green-lemma-right}
Für $a,b\in S$ und $u,v\in S^1$ mit $au=b$ und $bv=a$
sind die Abbildungen
\[
\begin{aligned}
L_a&\longrightarrow L_b,&x&\longmapsto xu,\\
L_b&\longrightarrow L_a,&y&\longmapsto yv
\end{aligned}
\]
Bijektionen mit den angegebenen Umkehrabbildungen. Für $x\in L_a$ gilt
$x\mathrel{\mathcal R}xu$; die erste Bijektion bildet $H_x$
bijektiv auf $H_{xu}$ ab. Insbesondere entsteht eine Bijektion
$H_a\longrightarrow H_b$.
\cite{harju,howie}
\end{theorem}

\begin{theorem}[Greens Lemma: Linksmultiplikation]
\label{b28guide:thm:green-lemma-left}
Für $a,b\in S$ und $u,v\in S^1$ mit $ua=b$ und $vb=a$
sind die Abbildungen
\[
\begin{aligned}
R_a&\longrightarrow R_b,&x&\longmapsto ux,\\
R_b&\longrightarrow R_a,&y&\longmapsto vy
\end{aligned}
\]
Bijektionen mit den angegebenen Umkehrabbildungen. Für $x\in R_a$ gilt
$x\mathrel{\mathcal L}ux$; die erste Bijektion bildet $H_x$
bijektiv auf $H_{ux}$ ab. Insbesondere entsteht eine Bijektion
$H_a\longrightarrow H_b$.
\cite{harju,howie}
\end{theorem}

\begin{theorem}[Raster und Gleichmächtigkeit innerhalb einer $\mathcal D$-Klasse]
\label{b28guide:thm:green-grid}
In jeder $\mathcal D$-Klasse $D$ schneidet jede darin enthaltene
$\mathcal R$-Klasse jede darin enthaltene $\mathcal L$-Klasse
in genau einer $\mathcal H$-Klasse. Alle diese Schnitte sind nichtleer.

Je zwei $\mathcal H$-Klassen in $D$ sind gleichmächtig.
Auch je zwei $\mathcal R$-Klassen in $D$ sind gleichmächtig,
und je zwei $\mathcal L$-Klassen in $D$ sind gleichmächtig.
Eine Gleichmächtigkeit zwischen einer $\mathcal R$-Klasse und einer
$\mathcal L$-Klasse wird damit nicht behauptet.
\cite{howie,green1951}
\end{theorem}

\section{Vollständige Alternative für $\mathcal H$-Klassen}

\begin{definition}[Untergruppe und maximale Untergruppe]
\label{b28guide:def:subgroup}
Eine \emph{Untergruppe von $S$} ist eine nichtleere,
produktabgeschlossene Teilmenge $G\subseteq S$ mit einem Element
$e\in G$, so daß
\[
(\forall g\in G)\;eg=ge=g,
\qquad
(\forall g\in G)(\exists h\in G)\;gh=hg=e.
\]
Neutralität und Inversen beziehen sich auf $G$; eine Gruppenstruktur
von ganz $S$ ist nicht gefordert. Eine Untergruppe heißt \emph{maximal},
wenn sie in keiner echt größeren Untergruppe von $S$ enthalten ist.
\end{definition}

\begin{theorem}[Greens vollständige $\mathcal H$-Klassen-Alternative]
\label{b28guide:thm:green-group}
Für jede $\mathcal H$-Klasse $H$ einer beliebigen Halbgruppe $S$
gilt genau eine der beiden Alternativen: $H$ ist eine maximale
Untergruppe von $S$, oder es gilt
\[
 H^2\cap H=\varnothing,\qquad H^2=\{ab:a,b\in H\}.
\]
Die Aussage erfaßt alle $\mathcal H$-Klassen; eine Gruppenklasse
wird nicht vorausgesetzt.
\cite{harju,howie}
\end{theorem}

\section{Schützenbergergruppen beliebiger Klassen}

\begin{definition}[Rechter Stabilisator und Schützenbergergruppe]
\label{b28guide:def:schutzenberger}
Für eine $\mathcal H$-Klasse $H$ setzen wir
\[
\begin{aligned}
\operatorname{Stab}_{R}(H)&=\{s\in S^1:Hs\subseteq H\},\\
\Gamma_R(H)&=\{\rho_s:s\in\operatorname{Stab}_{R}(H)\},
\end{aligned}
\]
wobei $\rho_s:H\longrightarrow H$, $h\longmapsto hs$.
Die Menge $\Gamma_R(H)$ mit der üblichen Komposition
$(\alpha\circ\beta)(h)=\alpha(\beta(h))$ heißt
\emph{rechte Schützenbergergruppe von $H$}.
\end{definition}

\begin{theorem}[Permutationssatz und einfache Transitivität]
\label{b28guide:thm:schutzenberger-action}
Für jede $\mathcal H$-Klasse $H$ und jedes $s\in S^1$ gilt
\[
Hs\subseteq H\iff Hs=H\iff Hs\cap H\ne\varnothing.
\]
Jede Abbildung $\rho_s$ mit $s\in\operatorname{Stab}_{R}(H)$ ist
bijektiv; $\Gamma_R(H)$ ist eine Gruppe von Permutationen von $H$.
Ihre Einheit ist $\operatorname{id}_H$, ihre Inversen sind
Umkehrabbildungen auf $H$. Multiplikative Inverse der Faktoren $s$
in $S^1$ werden nicht vorausgesetzt.
Für $s,t\in\operatorname{Stab}_{R}(H)$ gilt bei der festgelegten
Kompositionsrichtung $\rho_{st}=\rho_t\circ\rho_s$.
Die Wirkung durch Auswertung ist \emph{einfach transitiv}:
\[
(\forall a,b\in H)(\exists!\,\alpha\in\Gamma_R(H))\;
\alpha(a)=b.
\]
Für jedes $a\in H$ ist daher
$\Gamma_R(H)\longrightarrow H$, $\alpha\longmapsto\alpha(a)$,
eine Bijektion.
\cite{brookes2024,howie}
\end{theorem}

\begin{theorem}[Invarianz des Gruppentyps in einer $\mathcal D$-Klasse]
\label{b28guide:thm:schutzenberger-invariance}
Für $a,b\in S$ gilt
\[
a\mathrel{\mathcal D}b
\Longrightarrow
\Gamma_R(H_a)\cong\Gamma_R(H_b).
\]
Die Aussage betrifft sämtliche $\mathcal H$-Klassen in einer
$\mathcal D$-Klasse.
\cite{brookes2024,howie}
\end{theorem}
